% =================================================
% Учительская версия рабочего листа. Урок 2
% =================================================

\worksheettitle
  {Классы. Чтение и запись чисел}
  {Урок 2 из 3. Версия для учителя}

\begin{teacherbox}[Цели урока]
  К концу урока ученик делит запись числа на классы справа налево, называет
  классы и их разряды, читает многозначные числа, записывает их цифрами,
  сохраняет нулевые и неполные классы.
\end{teacherbox}

\begin{checkpointbox}[Входная диагностика: ответы]
  \begin{enumerate}[label=\textbf{\arabic*.},itemsep=0.35em]
    \item Цифра \(2\).
    \item \(70\,308=70\,000+300+8\).
    \item \(40\,702\).
    \item Семь цифр.
  \end{enumerate}
\end{checkpointbox}

\begin{teacherbox}[Как читать диагностику]
  Ошибки в заданиях 2--3 означают, что ученик ещё не удерживает нулевые
  разряды. До перехода к классам дайте ему одну-две записи в разрядной
  таблице. Ошибка в подсчёте цифр часто возникает из-за игнорирования нулей.
\end{teacherbox}

\section{Делим число на классы}

\[
  5308=5\,308,\qquad
  7204061=7\,204\,061,
\]
\[
  4805070204=4\,805\,070\,204.
\]

\splitclassesfigure

\begin{conclusionbox}
  Число делят на классы, начиная с правого края, по три цифры.
\end{conclusionbox}

\begin{teacherbox}[Методический акцент]
  Просите ученика закрыть часть числа слева и группировать запись именно
  справа. Деление слева особенно часто даёт ошибку в числах, количество цифр
  которых не кратно трём: например, \(72\,040\,61\).
\end{teacherbox}

\section{Названия классов}

I класс --- единицы; II --- тысячи; III --- миллионы; IV --- миллиарды.

\classtablefigure

Для числа \(7\,304\,052\,018\):
\begin{itemize}
  \item миллиарды --- \(7\);
  \item миллионы --- \(304\);
  \item тысячи --- группа \(052\), число внутри класса \(52\);
  \item единицы --- группа \(018\), число внутри класса \(18\).
\end{itemize}

\begin{teacherbox}[Группа \(052\) и число \(52\)]
  Полезно различать две формулировки: «как записан класс?» --- \(052\) и
  «сколько единиц класса?» --- \(52\). Начальные нули нужны для сохранения
  трёх позиций внутри непервого класса.
\end{teacherbox}

\section{Работаем с разрядной таблицей}

\classworksheetfigure

Для числа \(325\,041\,708\):
\begin{itemize}
  \item класс миллионов содержит \(325\);
  \item класс тысяч содержит \(41\) и записан группой \(041\);
  \item класс единиц содержит \(708\);
  \item цифра сотен тысяч равна \(0\);
  \item цифра десятков единиц равна \(0\).
\end{itemize}

\begin{teacherbox}[Ответ на вопрос о группе \(041\)]
  В каждом непервом слева классе должны быть видны все три разряда:
  сотни тысяч, десятки тысяч и единицы тысяч. Сотен тысяч нет, поэтому их
  место занимает ноль.
\end{teacherbox}

\section{Читаем многозначные числа}

\readingroutefigure

\begin{fillbox}[Ответы]
  \begin{enumerate}
    \item \(4\,218\) --- четыре тысячи двести восемнадцать.
    \item \(73\,005\) --- семьдесят три тысячи пять.
    \item \(6\,040\,019\) --- шесть миллионов сорок тысяч девятнадцать.
    \item \(502\,007\,030\) --- пятьсот два миллиона семь тысяч тридцать.
  \end{enumerate}
\end{fillbox}

\zeroclassfigure

В числе \(9\,000\,406\) нулевой класс --- класс тысяч. Его название не
произносят. Число читают: «девять миллионов четыреста шесть».

\begin{teacherbox}[Как проверить понимание]
  Дайте рядом \(9\,000\,406\) и \(9\,406\). Попросите назвать значение
  цифры \(9\): в первом числе это \(9\) миллионов, во втором --- \(9\) тысяч.
  Так ученик видит, почему группу \(000\) нельзя удалить.
\end{teacherbox}

\section{Записываем число цифрами}

\writingclassesfigure

\begin{fillbox}[Заполненная таблица]
  \renewcommand{\arraystretch}{1.5}
  \begin{tabularx}{\linewidth}{@{}X>{\centering\arraybackslash}p{0.16\linewidth}>{\centering\arraybackslash}p{0.16\linewidth}>{\centering\arraybackslash}p{0.16\linewidth}@{}}
    \toprule
    \rowcolor{TableHeader}
    Число словами & миллионы & тысячи & единицы \\
    \midrule
    сорок два миллиона семь тысяч девяносто & \(42\) & \(007\) & \(090\) \\
    триста пять миллионов восемьсот два & \(305\) & \(000\) & \(802\) \\
    восемнадцать миллионов сорок тысяч пять & \(18\) & \(040\) & \(005\) \\
    \bottomrule
  \end{tabularx}
\end{fillbox}

Получаем:
\[
  42\,007\,090,\qquad
  305\,000\,802,\qquad
  18\,040\,005.
\]

\begin{fillbox}[Запись цифрами]
  \begin{enumerate}
    \item шесть миллионов сорок два --- \(6\,000\,042\);
    \item семьдесят миллионов три тысячи восемь --- \(70\,003\,008\);
    \item девять миллиардов три миллиона двадцать тысяч четыре.\par
      Ответ: \(9\,003\,020\,004\).
  \end{enumerate}
\end{fillbox}

\begin{teacherbox}[Рабочий приём]
  На первых примерах разрешите подписывать над группами «млн | тыс | ед».
  Затем уберите подписи, но сохраните требование: каждый класс справа от
  первого записывается тремя цифрами.
\end{teacherbox}

\section{Исправляем ошибки}

\begin{warningbox}[Ошибка 1]
  Потерян промежуточный класс тысяч. Верно:
  \[
    \text{шесть миллионов сорок два}=6\,000\,042.
  \]
\end{warningbox}

\begin{warningbox}[Ошибка 2]
  В чтении числа \(53\,008\,004\) пропущены слова «восемь тысяч».
  Верно: «пятьдесят три миллиона восемь тысяч четыре».
\end{warningbox}

\begin{teacherbox}[Диагностическое различие]
  Если ученик пишет \(6\,042\), он теряет целый класс. Если пишет
  \(6\,000\,42\), он понимает наличие класса, но ещё не дополняет его до
  трёх разрядов. Эти ошибки требуют разных коррекционных вопросов.
\end{teacherbox}

\section{Самостоятельная проверка: ответы}

\begin{resultbox}
  \begin{enumerate}
    \item \(805070204=805\,070\,204\).
    \item \(6\) (группа \(006\)).
    \item Сорок миллионов пять тысяч восемнадцать.
    \item \(203\,070\,009\).
    \item \(4\,012\,000\,005\).
  \end{enumerate}
\end{resultbox}

\begin{checkpointbox}[Выходной билет: ответы]
  \begin{enumerate}
    \item \(014\).
    \item Три миллиона двести семь.
    \item \(50\,000\,008\).
  \end{enumerate}
\end{checkpointbox}

\begin{teacherbox}[Критерий готовности к уроку 3]
  Ученик готов к составлению чисел из цифр, если он:
  \begin{itemize}
    \item уверенно группирует запись справа;
    \item правильно читает числа с нулевыми классами;
    \item дополняет неполные классы слева нулями;
    \item выполняет не менее четырёх из пяти заданий самостоятельной проверки.
  \end{itemize}
\end{teacherbox}
